\begin{fixedbox}
\textbf{Downstream consequence-facing only.} This section is archival packaging downstream of Sections~1--9 and of the placement rules in Section~\ref{sec:consequence-pathways}. It records formalization posture, dependency routing, source-role discipline, and audit consequences of the already closed theorem chain. Nothing in this section adds a new proof burden to the foundational arc.
\end{fixedbox}

\section{Downstream audit consequences: formalization posture and archival architecture}
\label{sec:formalization-audit}

The substantive theorem chain is already closed before this section begins. What remains is downstream audit architecture: what depends on what, what is proved locally, what is imported, what Lean support is actually documented, what a hostile reader must deny to reject the paper, and which sources are permanently outside the live theorem arc.

\begin{importnote}[Audit sources used here]
This section uses only the documentation actually supplied by the declared source set: \fcas{} for the fixed-kernel formalization boundary and documented package facts; \trajpaper{} for the public theorem surface, theorem-family correspondences, and dependency spine of the trajectory-side Lean development; \noninstgodeldoc{} for the manuscript-specific axiom-free theorem-ladder audit; \standeqdoc{} for the audit-grade elimination of residual dual-track governance between admissibility and standing; \exhaustiondoc{} for typed boundary discipline, fixed-signature factorization against same-scope discriminators, and the capstone-facing necessity-layer witnesses; and the completed capstone itself as the theorem object whose dependencies are now being frozen. No broader Lean claim is taken from the latter two sources than those audit-facing closure roles.
\end{importnote}

\subsection{Downstream formalization posture}

\begin{definition}[Capstone formalization boundary]
The \emph{capstone formalization boundary} is the exact point at which the manuscript stops making formalization claims stronger than the supplied sources support. In the present manuscript that boundary is as follows:
\begin{enumerate}
  \item the kernel floor is \emph{not} claimed as a theorem generated inside a lower internal Lean layer;
  \item documented Lean surfaces are cited only where a source manuscript names public theorem wrappers, theorem families, package facts, or a manuscript-specific audit statement; and
  \item no capstone-wide merged Lean development is claimed unless a dedicated capstone audit exists in the declared source set.
\end{enumerate}
\end{definition}

\begin{theorem}[No-overclaim formalization rule]
\label{thm:no-overclaim-formalization}
The capstone may claim formalization support only in one of the following forms:
\begin{enumerate}
  \item \emph{theorem-level} support, when the declared sources document exact public theorem or wrapper names;
  \item \emph{theorem-family or module-level} support, when the declared sources document only a family-level or module-level correspondence; or
  \item \emph{manuscript-specific audit} support, when a source manuscript states that its own theorem chain has been encoded in Lean and its exported paper-facing theorems are axiom-free, but does not expose the full public name surface in the retrieved documentation.
\end{enumerate}
No fourth, stronger form is admissible here.
\end{theorem}

\begin{proof}
Theorem-level support is the strongest form explicitly documented in \fcas{} and \trajpaper{}. Theorem-family or module-level support is also explicitly documented there when the paper-level theorem compresses several lower-level Lean results or infrastructure layers. The manuscript-specific audit form is the strongest support documented in \noninstgodeldoc{} for the same-domain G\"odel non-instantiability family. A stronger capstone-wide claim would require a dedicated unified Lean audit, which is not present in the declared source set. So no fourth, stronger form is admissible.
\end{proof}

\begin{corollary}[Present capstone formalization posture]
The present capstone is auditable in two exact ways and in no stronger one:
\begin{enumerate}
  \item as an in-paper theorem chain whose proof burden is carried by the main sections and their declared source roles; and
  \item as a layered manuscript whose major theorem families can be linked to documented Lean surfaces at theorem, theorem-family, or manuscript-audit level where such documentation exists.
\end{enumerate}
It is not licit on the present evidence to claim one unified capstone Lean development with a complete exported theorem surface.
\end{corollary}

\begin{remark}[Formalization posture of the fixed floor]
Within the declared source set, the fixed floor is not treated as an omitted lower theorem waiting to be recovered by proof-assistant work. It is treated as a theorem-level boundary on what same-domain derivation could mean at all, while the downstream uniqueness, collapse, transport, irrecoverability, and non-instantiability families are exactly the parts for which theorem-level, theorem-family, or manuscript-audit Lean support is claimed. The formalization posture is therefore faithful to the proof order of the manuscript rather than weaker than it.
\end{remark}

\subsection{Downstream audit seal: axiom-necessity layer and closed-system closure}

\begin{definition}[Capstone axiom-necessity layer]
For audit purposes, collect the governing closure premises already used by Sections~\ref{sec:modeling-necessity}--\ref{sec:meta-closure} under the name \emph{capstone axiom-necessity layer}:
\begin{enumerate}
  \item evaluated-fragment bivalence together with fail-closed governance;
  \item finite same-scope compositional closure for candidate acts, histories, and continuations;
  \item lawful redescription neutrality together with stable reference;
  \item non-degeneracy of the evaluated reasoning regime;
  \item irreversibility and the ban on retroactive repair of standing-bearing lineage;
  \item no undeclared same-scope boundary selector, discriminator, or gate-multiplying authority;
  \item identity discipline for same-scope operator or representative copies, so that label-only duplication cannot create standing-bearing difference; and
  \item well-founded local denial-target discipline for hostile-reader audit and failure localization.
\end{enumerate}
These are not new starting axioms of the capstone. They are the audit-facing names for the governing premises whose necessity is exposed by the completed theorem chain and by the witness package consolidated in \exhaustiondoc{} \cite{MaleyExhaustion}.
\end{definition}

\begin{theorem}[Axiom-necessity layer for non-degenerate reasoning]
\label{thm:axiom-necessity-layer}
Each member of the capstone axiom-necessity layer is individually necessary for non-degenerate reasoning in the sense of this manuscript. More precisely:
\begin{enumerate}
  \item weakening bivalent fail-closed governance reopens intermediate or competing standing classifications and destroys same-scope gate uniqueness;
  \item weakening finite same-scope compositional closure reopens deferred admissibility and removes the forcing of pre-evaluation gate discipline;
  \item weakening lawful redescription neutrality reopens presentation-dependent standing and non-factorizing same-scope discriminators;
  \item weakening non-degeneracy collapses the regime into vacuous all-accept or all-reject form;
  \item weakening irreversibility reopens retroactive repair, reinterpretation, and same-lineage reset;
  \item weakening boundary closure reopens selectable competing gates, metadata-level authorizers, and undeclared same-scope scope drift;
  \item weakening identity discipline for operator or representative copies reopens label-dependent splitting beneath unchanged standing content; and
  \item weakening well-founded local denial-target discipline destroys unique minimal failure localization and with it the audit-level routing of objections.
\end{enumerate}
Hence none of these governing premises is dispensable if the capstone is to remain a non-degenerate same-scope closure theorem.
\end{theorem}

\begin{proof}
Sections~\ref{sec:modeling-necessity} and \ref{sec:standing-regime} already force evaluated-fragment bivalence, fail-closed admissibility, and unique same-scope classification; allowing intermediate or fail-open standing statuses would therefore reinstall same-side splits or competing gates beneath the same fixed scope. Finite same-scope compositional closure is required because the capstone's gate theorems and irrecoverability claims are about lawful reuse, continuation, and recombination; if those candidate composites need not be formable, deferred admissibility reappears and pre-evaluation gate discipline is no longer forced. Lawful redescription neutrality and stable reference are necessary because otherwise standing can vary by presentation alone, which is exactly the representative-level non-factorization blocked by Section~\ref{sec:standing-regime}. Non-degeneracy is necessary because without at least one live admissible/inadmissible split the closure theorem becomes vacuous.

Irreversibility is necessary because otherwise same-act repair, retrospective reinterpretation, or lineage reset reopen, contradicting the no-repair family and the same-domain meta-closure results. Boundary closure is necessary because every same-scope selector, extra discriminator, metadata-level authorizer, or silently imported gate performs the very authority multiplication ruled out by the AMetric boundary and by Theorem~\ref{thm:same-scope-factorization-discipline}. Identity discipline for operator or representative copies is necessary because otherwise label-only duplication or representative privileging can generate standing-bearing difference beneath unchanged quotient content, reopening the representative-smuggling routes blocked in Sections~\ref{sec:standing-regime} and \ref{sec:operator-exhaustion}. Finally, well-founded local denial-target discipline is necessary for audit because the hostile-reader architecture of this capstone requires a unique minimal place at which a denial must land; without that discipline, failure localization blurs and the audit map ceases to route objections to a determinate theorem-bearing burden. The witness package of \exhaustiondoc{} records the short countermodel form of each of these failures, while the capstone sections named above show that none of them can be relaxed without reopening a forbidden route.
\end{proof}

\begin{corollary}[Formal closed-system seal]
\label{cor:formal-closed-system-seal}
Relative to the declared source set, the capstone is closed in the formal same-scope sense: every internal objection must do exactly one of the following:
\begin{enumerate}
  \item deny one member of the capstone axiom-necessity layer;
  \item deny a named theorem whose proof depends on that member; or
  \item declare a scope change.
\end{enumerate}
No fourth same-scope internal route remains.
\end{corollary}

\begin{proof}
Theorem~\ref{thm:axiom-necessity-layer} removes the last ``dispensable premises'' objection by showing that each governing closure premise carries a corresponding failure mode when weakened. Section~\ref{sec:standing-regime} closes residual dual-track governance and same-scope discriminator escape by Theorems~\ref{thm:standing-admissibility-equivalence} and \ref{thm:same-scope-factorization-discipline}. Section~\ref{sec:meta-closure} then exhausts the same-domain meta-level rescue routes from above, especially in Theorems~\ref{thm:rescue-route-exhaustion} and \ref{thm:incompleteness-route-exhaustion}. So an objection that remains inside the same scope has no unclassified route left: it must identify the governing premise or theorem it denies. If it introduces new same-scope authority instead, the move is a scope change rather than an internal objection.
\end{proof}

\begin{remark}[What the necessity layer does and does not do]
The present subsection does not convert the capstone into a free-standing axiom-starting paper. It does something narrower: it makes explicit that the governing closure premises already used by the manuscript are individually necessary for non-degenerate same-scope reasoning, so that the remaining hostile-reader move is forced from vague dissatisfaction into named premise denial or scope change.
\end{remark}

\subsection{Downstream dependency packaging}

\begin{theorem}[Final dependency register theorem]
\label{thm:final-dependency-register}
As a downstream audit consequence of the completed closure theorem, the capstone has one strict dependency order and no hidden parallel proof route:
\begin{enumerate}
  \item Section~2 fixes the kernel floor and mechanization boundary.
  \item Section~3 forces state-transition form, identity persistence, irreversible commitment, and governance-level bivalence.
  \item Section~4 forces the AMetric boundary, relabeling-indifference collapse, no selector / no order / no deeper invariant, and singleton admissibility.
  \item Section~5 fixes gate existence, standing emergence, standing--admissibility equivalence, same-scope factorization against standing-relevant discriminators, no repairs / no generators / no carriers, standing conservation, the standing quotient, the unique admissible interior, and the compressed no-alternative-architecture theorem for the same-domain interface.
  \item Section~6 proves operator exhaustion and mathematics closure.
  \item Section~7 proves physics closure and descriptive non-generativity.
  \item Section~8 proves same-domain meta-closure and the non-instantiability of G\"odel-style escape routes.
  \item Section~9 proves two-spectrum exhaustion, closure parity, mirror classification, the prohibition package, and descriptive closure of admissible description.
  \item Section~10 places the downstream consequence manuscripts without allowing them to backflow into the already closed theorem chain.
  \item The present section adds no new foundational burden; it stabilizes the audit architecture of the completed chain, exposes the axiom-necessity layer for the governing closure premises, and states the formal closed-system corollary of the already completed theorem chain.
\end{enumerate}
No later section reopens an earlier burden without explicit revision of the declared source roles.
\end{theorem}

\begin{proof}
This is exactly the dependency discipline enforced by the capstone as rewritten. Each section depends only on the sections and declared sources activated before it, and Theorem~\ref{thm:no-downstream-backflow} already rules out application backflow from Section~\ref{sec:consequence-pathways}. So the dependency order is strict and non-circular.
\end{proof}

\begin{reviewbox}
\textbf{Audit architecture principle.} This section is archival, not generative. Its role is to make the theorem chain inspectable, not to supply missing proof content by packaging rhetoric.
\end{reviewbox}
